\documentclass[11pt]{article}
\usepackage{mathblatt}
% Probe 08.10.2026: „Die lange Seite“ in drei Durchgängen, Durchgang 3 (Prüfung), von Hand gesetzt
% nach aufgabenbank bau/bauregeln.md (Stand 08.10.2026) und mathe-nachhilfe plan.md § 5 „Richtung Durchgänge“.
% Prüfung: ohne Hilfe, ohne Päckchen, am Ende das Original ganz; gleiche Sorten, andere Varianten.
% Gemeinsame Präambel der drei Durchgänge, übernommen aus paket-pythagoras-2/2-blatt.tex (Probe 08.10.2026).
\makeatletter
\setlength{\mbpfbe}{0mm}% keine Punkte (Bauregel 6.4)
% eingerückt (Bauregel 4.7, 6.6): \gEin Einzug, \gSchrift eine Stufe kleiner
\newlength{\gEin}\setlength{\gEin}{0pt}
\let\gSchrift\relax
\renewcommand{\pfkreuzzeile}[1]{\par\noindent\hangindent3.2em\hangafter1\hspace*{1.6em}$\square$\ #1\par}
\newcommand{\gkreuz}[1]{$\square$\ #1\hspace{2em}}
% Bauregel 6.7: links, was man ansieht, rechts, was man tut; Spaltengrenze überall an derselben Stelle
\newsavebox{\gbox}\newlength{\gLinks}
\newcommand{\gzwei}[2]{\par\smallskip\noindent\sbox{\gbox}{#1}%
  \setlength{\gLinks}{\dimexpr0.5\textwidth-\gEin-\mbpfnr\relax}%
  \ifdim\wd\gbox>\dimexpr\gLinks-4mm\relax
    \usebox{\gbox}\par\smallskip\noindent\begin{minipage}{\linewidth}#2\end{minipage}%
  \else
    \begin{minipage}[t]{\gLinks}\vspace{0pt}\usebox{\gbox}\end{minipage}%
    \begin{minipage}[t]{\dimexpr\linewidth-\gLinks\relax}\vspace{0pt}\raggedright #2\end{minipage}%
  \fi\par}
\RenewDocumentEnvironment{pfaufg}{m m m}{%
  \begin{lrbox}{\mbaufgabebox}\begin{minipage}[t]{\linewidth}%
  \gSchrift\noindent\hspace*{\gEin}\mbpfrandmarke{#1}%
  \begin{minipage}[t]{\mbpfnr}\strut\textbf{#2}\end{minipage}%
  \begin{minipage}[t]{\dimexpr\linewidth-\gEin-\mbpfnr\relax}\raggedright\strut\ignorespaces}%
 {\end{minipage}%
  \end{minipage}\end{lrbox}%
  \setlength{\mbaufgabehoehe}{\dimexpr\ht\mbaufgabebox+\dp\mbaufgabebox\relax}%
  \par\addvspace{7pt}\Needspace*{\mbaufgabehoehe}%
  \noindent\usebox{\mbaufgabebox}\par\addvspace{7pt}}
\makeatother
% Rechenraum als Karo (Bauregel 6.9): n Rechenschritte = n mal 10 mm
\newcommand{\karo}[1]{\par\smallskip\noindent\begin{tikzpicture}
  \clip (0,0) rectangle ({\linewidth-1mm},{#1*10mm});
  \draw[black!16,line width=0.3pt,step=5mm] (0,0) grid ({\linewidth},{#1*10mm});
  \end{tikzpicture}\par}
\newcommand{\kk}[1]{$\square$\,#1\hspace{0.9em}}
\newcommand{\linie}[1]{\,{\color{black!45}\rule[-1pt]{#1}{0.4pt}}\,}
% Zeile am Ende jedes Durchgangs: Originale klein grau, darüber; dann Vorgänger/Nachfolger (Bauregel 3.4, 6.5)
\newcommand{\blattende}[2]{\par\vfill\noindent\begin{minipage}{\linewidth}{\scriptsize\color{mbgrau}Originale: #1\par}\smallskip
{\footnotesize #2\par}\end{minipage}}
\newcommand{\pkt}{\textperiodcentered{}}

\newcommand{\kennung}{H8Q}

\begin{document}
\pfheftstil
\fancyfoot[C]{\parbox[t]{\textwidth}{\mbfhbox\par\vspace{1.5mm}{\footnotesize\color{mbgrau}{\scriptsize \kennung}\hfill\thepage}}}
\pfheftkopf{Die lange Seite \pkt{} Hypotenuse}{P10 \pkt{} 3}

% 1 Satz in Worten (Kette 2 s7 v3)
\begin{pfaufg}{}{1.}{}
Kreuze die Aussage an, die in jedem rechtwinkligen Dreieck gilt.
\pfkreuzzeile{Die Katheten liegen am rechten Winkel; ihre Quadrate ergeben zusammen das Quadrat der dritten Seite.}
\pfkreuzzeile{Die Hypotenuse liegt am rechten Winkel; ihr Quadrat ist die Summe der anderen zwei.}
\pfkreuzzeile{Das Quadrat der längsten Seite ist kleiner als die zwei anderen Quadrate zusammen.}
\fusshilfe{\mbox{1:~die erste Aussage}}
\end{pfaufg}

% 2 Gleichung aufstellen (Kette 2 s6 v2), ohne Hilfe; Hypotenuse steht senkrecht links
\begin{pfaufg}{}{2.}{}
Schreibe für dieses Dreieck die Gleichung nach dem Satz des Pythagoras auf.
\gzwei{\begin{tikzpicture}[line width=0.7pt,scale=0.55]\coordinate (P) at (0,2.6);\coordinate (Q) at (0,0);\coordinate (R) at (1.2,1.8);\draw (P) -- (Q) -- (R) -- cycle;\mbrwmarke{(R)}{(P)}{(Q)}
\node[font=\small,anchor=east] at (-0.08,1.3) {$f$};\node[font=\small,anchor=south west] at (0.62,2.22) {$g$};\node[font=\small,anchor=north west] at (0.66,0.86) {$e$};\end{tikzpicture}}{%
\vspace{4mm}\linie{50mm}}
\fusshilfe{\mbox{2:~$g^2 + e^2 = f^2$}}
\end{pfaufg}

% 3 Formel mit Wurzel (Kette 2 s8 v3); rechter Winkel oben, Hypotenuse unten
\begin{pfaufg}{}{3.}{}
Kreuze die Formel an, mit der man $r$ berechnet.
\gzwei{\begin{tikzpicture}[line width=0.7pt,scale=0.7]\coordinate (A) at (0,0);\coordinate (B) at (3.2,0);\coordinate (C) at (0.8,1.386);\draw (A) -- (B) -- (C) -- cycle;\mbrwmarke{(C)}{(A)}{(B)}
\node[font=\small,anchor=south east] at (0.38,0.72) {$p$};\node[font=\small,anchor=south west] at (2.05,0.72) {$q$};\node[font=\small,anchor=north] at (1.6,-0.08) {$r$};\end{tikzpicture}}{%
\pfkreuzzeile{$r = \sqrt{p^2 - q^2}$\hspace{2.4em}$\square$\ $r = \sqrt{q^2 - p^2}$}
\pfkreuzzeile{$r = q^2 + p^2$\hspace{3.7em}$\square$\ $r = \sqrt{p^2 + q^2}$}}
\fusshilfe{\mbox{3:~$r = \sqrt{p^2 + q^2}$}}
\end{pfaufg}

% 4 Hypotenuse berechnen, eine Aufgabe statt Päckchen (Kette 2 s4 v3: Kommazahlen und Runden)
\begin{pfaufg}{}{4.}{}
Die Katheten sind $3{,}6$~m und $4{,}1$~m lang. Berechne die Hypotenuse $c$ auf eine Stelle nach dem Komma.
\karo{3}\par\smallskip\hfill $c \approx$\linie{18mm}m
\fusshilfe{\mbox{4:~$c = \sqrt{29{,}77} \approx 5{,}5$ m}}
\end{pfaufg}

% 5 Dreieck in anderer Lage (Kette 2 s2 v2), gekippt
\begin{pfaufg}{}{5.}{}
Im Dreieck ist eine Seite mit ? markiert. Wie lang ist diese Seite?
\gzwei{\begin{tikzpicture}[line width=0.7pt,scale=0.9]\coordinate (R) at (1.6,0);\coordinate (L) at (0.049,0.564);\coordinate (M) at (2.558,2.631);\draw (R) -- (L) -- (M) -- cycle;\mbrwmarke{(R)}{(L)}{(M)}
\node[font=\small,anchor=north east] at (0.82,0.26) {$33$ cm};\node[font=\small,anchor=west] at (2.15,1.25) {$56$ cm};\node[font=\small,anchor=south east] at (1.28,1.64) {?};\end{tikzpicture}}{%
\karo{3}\par\smallskip\hfill ? $=$\linie{18mm}cm}
\fusshilfe{\mbox{5:~65 cm}}
\end{pfaufg}

% 6 Diagonale im Rechteck (Kette 2 s9 v3), ohne Skizze
\begin{pfaufg}{}{6.}{}
Ein Rechteck ist $9$~cm lang und $7$~cm breit. Berechne die Diagonale auf eine Stelle nach dem Komma.
\karo{3}\par\smallskip\hfill Diagonale $\approx$\linie{18mm}cm
\fusshilfe{\mbox{6:~$\sqrt{130} \approx 11{,}4$ cm}}
\end{pfaufg}

% 7 Sache mit Skizze, Prüfungshöhe (Kette 2 s15 v4, eigene Variante zu 2025 · 4 · a)
\begin{pfaufg}{}{7.}{}
Ein Steg führt hinunter zu einem Boot. Berechne seine Länge auf eine Stelle nach dem Komma.
\gzwei{\begin{tikzpicture}[line width=0.7pt,scale=0.62]\coordinate (P) at (0,0);\coordinate (Q) at (5.95,0);\coordinate (U) at (0,0.6);\draw[line width=0.4pt] (P) -- (Q);\draw (P) -- (U);\draw[line width=1.4pt] (U) -- (Q);\mbrwmarke{(P)}{(Q)}{(U)}
\node[font=\small,anchor=north] at (2.97,-0.08) {$350$ cm};\node[font=\small,anchor=east] at (-0.08,0.3) {$22$ cm};\node[font=\small,anchor=south] at (3.1,0.32) {Steg};
\node[font=\scriptsize,mbgrau,anchor=north west] at (0,-0.6) {nicht maßstabsgerecht};\end{tikzpicture}}{%
\karo{3}\par\smallskip\hfill Steg $\approx$\linie{18mm}cm}
\fusshilfe{\mbox{7:~$\sqrt{122\,984} \approx 350{,}7$ cm}}
\end{pfaufg}

% 8 Überstand (Kette 2 s11 v3), ohne Hilfe
\begin{pfaufg}{}{8.}{}
Ein Dachsparren ragt unten über die Hauswand hinaus. Wie lang ist der ganze Sparren?
\gzwei{\begin{tikzpicture}[line width=0.7pt,scale=0.62]
\fill[black!10] (-0.35,-2.0) rectangle (0,0);\draw (0,-2.0) -- (0,0);
\draw[line width=1.6pt] (-0.529,-0.282) -- (4.5,2.4);
\draw[dashed,line width=0.4pt] (0,0) -- (4.5,0) -- (4.5,2.4);
\mbrwmarke{(4.5,0)}{(0,0)}{(4.5,2.4)}
\node[font=\small,anchor=north] at (2.25,-0.08) {$4{,}5$ m};
\node[font=\small,anchor=west] at (4.58,1.2) {$2{,}4$ m};
\node[font=\small,anchor=south east] at (-0.2,-0.1) {$0{,}6$ m};
\node[font=\scriptsize,mbgrau,anchor=north west] at (0.2,-0.7) {Hauswand links};
\end{tikzpicture}}{%
\karo{4}\par\smallskip\hfill Sparren:\linie{18mm}m}
\fusshilfe{\mbox{8:~$5{,}1 + 0{,}6 = 5{,}7$ m}}
\end{pfaufg}

% 9 Fehler finden (Kette 5 s1 v3): begründen ohne Rechnung
\begin{pfaufg}{}{9.}{}
Tim hat vier Hypotenusen berechnet. Welche Ergebnisse können nicht stimmen? Begründe, ohne genau zu rechnen.
\gzwei{\begin{minipage}{62mm}\small
(1)\ $12$ m und $35$ m:\quad $c = 37$ m\par
(2)\ $11$ cm und $60$ cm:\quad $c = 49$ cm\par
(3)\ $18$ m und $80$ m:\quad $c = 98$ m\par
(4)\ $13$ cm und $84$ cm:\quad $c = 85$ cm\end{minipage}}{\karo{3}}
\fusshilfe{\mbox{9:~(2): kürzer als eine Kathete; (3): so lang wie beide Katheten zusammen}}
\end{pfaufg}

% 10 am Ende das Original ganz (2020 · 7 · a), mit Punkt Q aus der Abbildung (Q braucht erst 7 b)
\begin{pfaufg}{P10 ’20}{10.}{}
Im Dreieck $ABC$ ist bei $C$ ein rechter Winkel. Die Strecke $\overline{BC}$ ist $12$~m lang, die Strecke $\overline{AC}$ $34$~m. Berechne die Länge der Strecke $\overline{AB}$.
\gzwei{\begin{tikzpicture}[line width=0.7pt]\coordinate (A) at (0,0);\coordinate (B) at (1.2,3.4);\coordinate (C) at (0,3.4);\coordinate (Q) at (0,2.7);\draw (A) -- (B) -- (C) -- cycle;\draw[dashed] (Q) -- (B);\mbrwmarke{(C)}{(A)}{(B)}\fill (Q) circle (1.4pt);
\node[font=\small] at (-0.05,-0.3) {$A$};\node[font=\small] at (1.37,3.64) {$B$};\node[font=\small] at (-0.1,3.68) {$C$};\node[font=\small,anchor=east] at (-0.06,2.7) {$Q$};
\node[font=\scriptsize,mbgrau,anchor=north west] at (0,-0.5) {nicht maßstabsgerecht};\end{tikzpicture}}{%
\karo{3}}
\fusshilfe{\mbox{10:~$\overline{AB} \approx 36{,}1$ m}}
\end{pfaufg}

\blattende{2020 \pkt{} 7 \pkt{} a}{Vorher: Die lange Seite \pkt{} P10 \pkt{} 2\quad\textbullet\quad Weiter: Eine kurze Seite \pkt{} Kathete}
\end{document}
